\documentclass[12pt,a4paper]{report}
\usepackage[utf8]{inputenc}
\usepackage{graphicx}
\usepackage{geometry}
\usepackage{titlesec}
\usepackage{tocloft}
\usepackage{fancyhdr}
\usepackage{setspace}
\usepackage{svg}

% Page Geometry
\geometry{top=1in, bottom=1in, left=1.25in, right=1in}

% Chapter Title Formatting
\titleformat{\chapter}[display]
  {\normalfont\LARGE\bfseries\centering}{\chaptertitlename\ \thechapter}{10pt}{\LARGE}
\titlespacing*{\chapter}{0pt}{-20pt}{20pt}

\begin{document}

% -------------------------------------------------------------------
% TITLE PAGE
% -------------------------------------------------------------------
\begin{titlepage}
\begin{center}
% Placeholder for CUSAT Logo
\begin{figure}[h]
    \centering
    \includesvg[width=4cm]{cusat.svg} 
\end{figure}

\vfill % --- FLEXIBLE GAP ---

\Large \textbf{PROJECT REPORT (PHASE 1)}\\[0.3cm]
\large \textbf{on}\\[0.3cm]
\LARGE \textbf{CARDIOKEY: ANALOG ECG SIGNAL ACQUISITION AND BIOMETRICS}

\vfill % --- FLEXIBLE GAP ---

\large \textit{Submitted by}\\[0.4cm]
{\large
\textbf{AARDRA S V – 24101868}\\
\textbf{AAROMAL A – 24101880}\\
\textbf{ABDULLA HILAL V K – 24101000}\\
\textbf{ABHIJITH R S – 24102771}}

\vfill % --- FLEXIBLE GAP ---

\textit{in partial fulfillment of requirement for the award of the}\\[0.2cm]
\textit{degree of}\\[0.4cm]
\Large \textbf{BACHELOR OF TECHNOLOGY}\\[0.2cm]
\large \textit{in}\\[0.2cm]
\Large \textbf{ELECTRONICS AND COMMUNICATION}

\vfill % --- FLEXIBLE GAP ---

\large \textbf{DIVISION OF ELECTRONICS ENGINEERING}\\
\textbf{SCHOOL OF ENGINEERING}\\
\textbf{COCHIN UNIVERSITY OF SCIENCE AND TECHNOLOGY}\\
\textbf{KOCHI--682022}\\[0.4cm]
\textbf{OCTOBER 2025}

\end{center}
\end{titlepage}
% -------------------------------------------------------------------
% CERTIFICATE PAGE
% -------------------------------------------------------------------
\newpage
\pagenumbering{roman}
\begin{center}
\large \textbf{DIVISION OF ELECTRONICS ENGINEERING}\\
\textbf{SCHOOL OF ENGINEERING}\\
\textbf{COCHIN UNIVERSITY OF SCIENCE AND TECHNOLOGY}\\
\textbf{KOCHI-682022}\\[1.5cm]

\Large \textbf{CERTIFICATE}\\[1cm]
\end{center}

\begin{spacing}{1.5}
\noindent \large \textit{Certified that the project report (phase 1) entitled \textbf{“CARDIOKEY: ANALOG ECG SIGNAL ACQUISITION AND BIOMETRICS”} is a bonafide work of \textbf{AARDRA S V (24101868)}, \textbf{AAROMAL A (24101880)}, \textbf{ABDULLA HILAL V K (24101000)}, and \textbf{ABHIJITH R S (24102771)}, towards the partial fulfilment for the award of the degree of B. Tech in Electronics and Communication of Cochin University of Science and Technology, Kochi-682022.}
\end{spacing}

\vspace{3cm}
\noindent \textbf{Project Guide} \hfill \textbf{Head of the Division}\\[0.5cm]
\noindent Ms. Reshmi R. Das \hfill Dr. Deepa Shankar

% -------------------------------------------------------------------
% ACKNOWLEDGEMENT
% -------------------------------------------------------------------
\newpage
\begin{center}
\Large \textbf{ACKNOWLEDGEMENT}\\[1cm]
\end{center}

\begin{spacing}{1.5}
\noindent We would like to express our sincere gratitude to the Cochin University of Science and Technology (CUSAT) and the Department of Electronics and Communication Engineering for providing us with the opportunity and resources to undertake this project titled “CardioKey: Analog ECG Signal Acquisition and Biometrics.”

\noindent We extend our heartfelt thanks to Dr. Deepa Shankar, Head of the Department of Electronics and Communication Engineering, for her support, encouragement, and for providing the necessary facilities and a motivating environment to carry out this work successfully.

\noindent We are deeply grateful to our Project Coordinators, Ms. Mrithula S and Ms. Babitha R Jose, for their continuous guidance, timely advice, and valuable coordination throughout the course of this project.

\noindent Our profound gratitude goes to our Project Guide, Ms. Reshmi R. Das, for her constant supervision, insightful suggestions, and technical guidance. Her encouragement and expertise played a vital role in shaping this project to its final form.

\noindent We also wish to express our appreciation to the laboratory faculty of the Department for their kind assistance and cooperation during the practical phases of the work.

\noindent Finally, we extend our sincere thanks to all faculty members of the Department for their valuable feedback during project reviews, and to our families and friends for their unwavering support and encouragement throughout this endeavor.
\end{spacing}

% -------------------------------------------------------------------
% ABSTRACT
% -------------------------------------------------------------------
\newpage
\begin{center}
\Large \textbf{ABSTRACT}\\[1cm]
\end{center}

\begin{spacing}{1.5}
\noindent Traditional biometric identification methods, such as fingerprint and facial recognition, are increasingly vulnerable to spoofing, fabrication, and external modification. The Electrocardiogram (ECG) presents a highly secure alternative as a biometric identifier due to its inherent uniqueness and the critical advantage of liveness detection. The core of the "CardioKey" project relies on the distinct anatomical and physiological characteristics of an individual's heart, which dictate the specific morphology of the P, QRS, and T waves. 

\noindent This project proposes a comprehensive system for analog ECG signal acquisition and biometric classification. The system utilizes a precision analog front-end, based on the AD620 instrumentation amplifier and an AD705J Right Leg Drive circuit, to acquire low-amplitude, Lead I heart signals while rejecting 60Hz mains hum and common-mode noise. The acquired analog signal is digitized using a 16-bit ADS1115 ADC and processed by a microcontroller. Digital Signal Processing (DSP) techniques, including wavelet drift correction and Butterworth filters, are applied to extract a clean P-QRS-T complex. Feature space reduction is performed via Principal Component Analysis (PCA), compressing the 250-sample heartbeat window into 30 principal components. Finally, classification algorithms such as Linear Discriminant Analysis (LDA) and Support Vector Machines (SVM) are utilized to identify the user. Experimental validation from foundational databases demonstrates an identification accuracy of up to 96\% under resting conditions, proving that the ECG is a viable, universal, and unique biometric trait for practical access control systems.
\end{spacing}

% -------------------------------------------------------------------
% TABLE OF CONTENTS
% -------------------------------------------------------------------
\newpage
\tableofcontents

% -------------------------------------------------------------------
% CHAPTERS
% -------------------------------------------------------------------
\newpage
\pagenumbering{arabic}

\chapter{INTRODUCTION}
As a core technology in the field of information security, human biometric recognition has become the focus of researchers' attention over the past few years. Traditional biometric methods, such as fingerprints and Face ID, are vulnerable to spoofing, latex replication, and photograph modification. Passwords and tokens can also be easily stolen or lost. 

To strengthen security, the electrocardiogram (ECG) has emerged as a powerful biometric identifier, guaranteeing reliability because it is strictly generated by living bodies. The ECG waveform possesses four fundamental characteristics required for biometric identification: universality, uniqueness, stability, and measurability. Every living human continuously generates ECG signals, and differences between individuals are driven by highly unique factors like heart location, geometric shape, and chest structure. 

The "CardioKey" project aims to build an analog implementation for ECG biometric identification. Muscle contraction during the cardiac cycle causes depolarization, which is detected as voltage changes on the skin's surface. By capturing these changes through a precision analog front-end and processing the unique morphology of the P, QRS, and T waves, the system can securely identify an individual.

\chapter{LITERATURE REVIEW}
Recent studies have established the efficacy of using ECG signals for human identification. Biel et al. extracted 30 time and amplitude features, utilizing Soft Independent Modeling of Class Analogy (SIMCA) to achieve a 98\% identification accuracy among 20 subjects. Shen et al. proposed a system based on single-lead ECG, securing 100\% correct verification by extracting 7 time-domain features from the QRS complex and employing decision-based neural networks.

Lugovaya and Nemirko investigated identification using the ECG-ID database, which involved 90 volunteers and 310 single-lead (Lead I) recordings. Their methodology extracted 250-sample P-QRS-T fragments, reduced dimensions via Principal Component Analysis (PCA), and utilized Linear Discriminant Analysis (LDA) to achieve a 96\% identification rate.

However, the robustness of ECG biometrics drops significantly under non-resting conditions. Cui et al. highlighted the "Rest-Ex" problem, observing that algorithms trained on resting data collapsed to a 2.7\% accuracy when tested on post-exercise data due to heart rate variability. By applying Kullback-Leibler (KL) divergence for feature selection, they improved the post-exercise accuracy to 61.4\%. 

\chapter{PROBLEM IDENTIFICATION AND OBJECTIVE}
The key problems observed in traditional access control and standard biometrics include:
\begin{itemize}
    \item Vulnerability to fabrication, spoofing attacks, and theft of credentials.
    \item Lack of inherent "liveness detection" in fingerprint or facial recognition.
    \item The "Rest-Ex" problem in ECG biometrics, where an elevated heart rate significantly distorts the baseline fragment of the ECG and attenuates the R-wave amplitude, causing identification algorithms to fail.
\end{itemize}

\section{Objectives}
The primary goal of this project is to design a robust analog acquisition and digital classification system for ECG biometrics. Specific objectives include:
\begin{enumerate}
    \item To design a precision analog front-end utilizing an instrumentation amplifier (AD620) to capture low-amplitude Lead I heart signals.
    \item To actively reject 60Hz mains hum and common-mode noise using a Right Leg Drive feedback circuit.
    \item To digitize the acquired signal at 500 Hz with 16-bit precision using an ADS1115 ADC.
    \item To implement digital preprocessing algorithms for baseline wander correction and high-frequency noise removal.
    \item To extract the informative P-QRS-T fragment consisting of 250 samples per heartbeat.
    \item To reduce feature space dimensionality via Principal Component Analysis (PCA) and successfully classify users via Support Vector Machines (SVM) or Linear Discriminant Analysis (LDA).
\end{enumerate}

\chapter{WHY WE CHOSE THIS PROJECT}
The team identified a growing need for highly secure, non-forgeable biometric systems in physical access control. We chose this project because ECG biometrics inherently provide "Liveness Detection," making it vastly superior to easily spoofed static images or silicone fingerprints. 

Furthermore, "CardioKey" represents a comprehensive multidisciplinary challenge. It bridges the gap between precision analog circuit design—requiring careful noise filtration of signals on the scale of millivolts—and advanced digital signal processing and machine learning. The use of a single-lead (Lead I) setup closely imitates the practical scenario of a user simply touching a biometric lock with their hands, making it a highly viable real-world application.

\chapter{PHASE 1 – ANALOG CIRCUIT DESIGN AND TESTING}
Phase 1 focused on the hardware acquisition of the raw analog ECG signal. Because the electrical signals from the heart are on the order of a few millivolts, a high-gain differential amplifier is required. 

The system utilizes an AD620 instrumentation amplifier due to its high common-mode rejection ratio and low input voltage noise ($9 nV/\sqrt{Hz}$). The gain is programmed accurately using a single external resistor, achieving a total circuit gain of 1001 ($1V/mV$ of heart signal).

To mitigate the immense 60Hz mains noise absorbed by the human body, an AD705J operational amplifier is configured as a Right Leg Drive (RLD) circuit. This circuit takes the common-mode voltage, inverts it, and feeds it back into the patient's right leg to actively cancel interference. 

\begin{figure}[htbp]
    \centering
    \rule{10cm}{6cm} % DUMMY PLACEHOLDER - INSERT PATH TO AD620 SCHEMATIC HERE
    % \includegraphics[width=0.8\textwidth]{path/to/AD620_schematic.png}
    \caption{Medical ECG Monitor Circuit utilizing AD620 and AD705J Right Leg Drive.}
\end{figure}

Testing on a breadboard confirmed the removal of high-frequency muscle noise via a passive low-pass filter and the elimination of baseline wander via a $0.03Hz$ high-pass filter. 

\chapter{SYSTEM DESIGN AND ARCHITECTURE}
The overall system architecture follows a classical biometric identification scheme, progressing strictly from hardware acquisition to software classification.

\section{Subsystems}
\begin{enumerate}
    \item \textbf{Data Acquisition:} Analog signals are gathered via Lead I (Left Arm - Right Arm) using limb clamp electrodes, mimicking a user interacting with a door handle.
    \item \textbf{Data Preprocessing:} The raw signal is cleaned using a 4th-order Butterworth bandpass filter (0.5 Hz to 40 Hz) and Wavelet Drift Correction.
    \item \textbf{Initial Feature Space Formation:} The continuous signal is segmented. The Pan-Tompkins algorithm detects R-peaks, and a 250-sample P-QRS-T window (80 samples before the peak, 170 after) is extracted.
    \item \textbf{Feature Space Reduction:} Raw 250-point data is compressed into 30 principal components via PCA.
    \item \textbf{Classification:} The 30-point signature is compared against a database using an SVM or LDA to determine identity.
\end{enumerate}

\begin{figure}[htbp]
    \centering
    \rule{10cm}{6cm} % DUMMY PLACEHOLDER - INSERT PATH TO SYSTEM STRUCTURE HERE
    % \includegraphics[width=0.8\textwidth]{path/to/system_structure.png}
    \caption{Detailed Identification System Structure.}
\end{figure}

\chapter{SENSORS AND EQUIPMENT OVERVIEW}
\section{Electrodes and Analog Front-End (AD8232 / AD620)}
The primary sensor interface relies on a 3-lead electrode system. Left Arm (LA) and Right Arm (RA) act as differential inputs to the AD620 instrumentation amplifier, while the Right Leg (RL) serves as the ground reference for common-mode noise cancellation. Alternatively, the fully integrated AD8232 heart monitor module can be utilized for simplified prototyping.

\section{Analog-to-Digital Converter (ADS1115)}
To satisfy the stringent resolution required for biometric feature extraction, the analog output is digitized using an ADS1115 module. This provides 16-bit precision and is configured to sample the waveform at 500 Hz. 

\section{Microcontroller Unit}
An ESP32 or Raspberry Pi acts as the central processing unit, reading the I2C data from the ADS1115 and handling the immediate DSP operations and matrix math required for classification.

\chapter{SIGNAL PROCESSING AND FILTERING}
Because the raw ECG signal is easily corrupted by baseline wander (from breathing) and high-frequency noise (from muscle movement), rigorous digital filtering is required. 

The pre-processing pipeline includes:
\begin{itemize}
    \item \textbf{Wavelet Drift Correction:} Multilevel 1D wavelet analysis (using biorthogonal wavelets at level 9) extracts and subtracts the drifting baseline. 
    \item \textbf{Adaptive Bandstop Filter:} Precisely targets and suppresses 60 Hz power-line noise without heavily distorting the QRS complex.
    \item \textbf{Butterworth Lowpass Filter:} Applied with a cut-off at 40 Hz to remove residual high-frequency muscle artifacts. 
\end{itemize}

\begin{figure}[htbp]
    \centering
    \rule{10cm}{6cm} % DUMMY PLACEHOLDER - INSERT PATH TO FILTERING RESULTS HERE
    % \includegraphics[width=0.8\textwidth]{path/to/filtering_results.png}
    \caption{Raw vs. Filtered ECG Signal demonstrating noise suppression.}
\end{figure}

\chapter{FEATURE EXTRACTION AND REDUCTION}
Initial feature space formation requires chopping the continuous ECG signal into discrete heartbeat segments. The prominent R-peak is identified utilizing the Hamilton Segmenter or Pan-Tompkins algorithm. 

Because the P-wave contains significant individual biometric data, the algorithm extracts a complete P-QRS-T fragment comprising exactly 250 samples (0.5 seconds). The fragment is aligned using the R-peak as a reference point.

To make classification computationally feasible, dimensionality reduction is mandatory. Principal Component Analysis (PCA) maps the 250-dimensional correlated feature space into an uncorrelated basis, extracting only the top 30 principal components representing the user's "signature".

\chapter{SOFTWARE IMPLEMENTATION}
The software backend, written in Python or MATLAB, orchestrates the data pipeline. 

\section{Pseudo-code Structure:}
\begin{enumerate}
    \item Read raw digitized ECG array via serial or I2C.
    \item Apply SciPy Butterworth filters and Wavelet transforms for noise reduction.
    \item Execute R-peak detection thresholds to calculate Heart Rate.
    \item Slice arrays into 250-sample P-QRS-T fragments.
    \item Feed fragments into the `sklearn.decomposition.PCA` module.
    \item Pass the resulting 30 vectors into the `sklearn.svm.SVC` (Support Vector Machine) model to verify user identity.
\end{enumerate}

\begin{figure}[htbp]
    \centering
    \rule{10cm}{6cm} % DUMMY PLACEHOLDER - INSERT PATH TO PQRST EXTRACTION HERE
    % \includegraphics[width=0.8\textwidth]{path/to/pqrst_extraction.png}
    \caption{Examples of R peak detection and P-QRS-T fragment extraction.}
\end{figure}

\chapter{RESULTS AND DISCUSSION}
Experimental evaluations based on the PhysioNet ECG-ID database demonstrated high viability for the system. Using Lead I data recorded at 500Hz with 12-bit precision, the identification system achieved an accuracy rate of 96\% under controlled resting conditions. 

However, testing revealed significant vulnerabilities to the "Rest-Ex" problem. When the classifier was trained on resting data and tested against post-exercise data (elevated heart rate), accuracy plummeted to 2.7\% using standard SVM techniques. By implementing Kullback-Leibler (KL) divergence for advanced feature selection—targeting only features insensitive to heart rate spikes—the post-exercise accuracy was recovered to 61.4\%. 

These results confirm that while ECG is a highly secure biometric at rest, advanced length normalization algorithms are critical for active environments.

\chapter{KEY FEATURES AND ADVANTAGES}
\begin{itemize}
    \item \textbf{Liveness Detection:} Inherently rejects fake fingerprints or photographs, as it requires a beating human heart.
    \item \textbf{Universality and Stability:} Every living human generates continuous, stable waveform structures determined by chest geometry.
    \item \textbf{Continuous Authentication:} Can continuously verify a user's identity as long as they are holding the handle/sensor, preventing mid-session hijacking.
    \item \textbf{High Rest Accuracy:} Proven 96\% identification accuracy using a simple single-lead setup.
\end{itemize}

\chapter{APPLICATIONS AND FUTURE SCOPE}
\section{Applications}
\begin{itemize}
    \item High-security physical access control systems (door locks, secure vaults).
    \item Continuous biometric authentication for steering wheels in smart vehicles.
    \item Patient identification and record-matching in remote medical monitoring.
\end{itemize}

\section{Future Scope}
\begin{itemize}
    \item \textbf{Overcoming HRV:} Implement advanced dynamic scaling of the QT interval to solve the Rest-Ex problem entirely.
    \item \textbf{Multi-modal Biometrics:} Combine the 1D ECG signal with fingerprints to create an un-spoofable dual-factor hardware lock.
    \item \textbf{Dry Electrodes:} Transition from sticky Ag/AgCl patches to stainless steel or copper touch-pads for seamless user interaction.
\end{itemize}

\chapter{COST ESTIMATION}
Table 14.1 details the estimated prototyping cost for the CardioKey hardware acquisition module.

\begin{table}[h]
\centering
\caption{Cost Estimation of Components}
\vspace{0.3cm}
\begin{tabular}{|l|c|r|}
\hline
\textbf{Item} & \textbf{Quantity} & \textbf{Approx. Cost (INR)} \\
\hline
Microcontroller (ESP32 / Raspberry Pi) & 1 & 1,500 \\
AD8232 ECG Sensor / AD620 ICs & 1 & 1,200 \\
ADS1115 16-bit ADC Module & 1 & 800 \\
Ag/AgCl Sticky Electrodes & 10 & 300 \\
Electrode Lead Cables (3.5mm jack) & 1 & 400 \\
Passive Components (Resistors, Capacitors) & - & 300 \\
Breadboard, Wiring, Misc. & - & 500 \\
\hline
\textbf{Total Estimated Cost} & & \textbf{5,000} \\
\hline
\end{tabular}
\end{table}

\chapter{CONCLUSION}
The "CardioKey" project successfully demonstrates the design and feasibility of using the Electrocardiogram (ECG) as a unique biometric identifier. By engineering a precise analog acquisition circuit utilizing the AD620, common noise is heavily attenuated, isolating the crucial P-QRS-T waveforms. 

The transition to digital signal processing allows for robust feature space reduction via PCA and accurate classification using machine learning models. While heart rate variability remains a persistent challenge for active users, the system's 96\% baseline accuracy proves that integrating analog biosignal acquisition with advanced digital classifiers creates a powerful, spoof-proof security solution.

\newpage
\begin{thebibliography}{99}
\bibitem{1} CardioKey (ECGID) - Zerothppt.pdf. Analog ECG Signal Acquisition \& Biometrics.
\bibitem{2} Lugovaya T.S. Biometric human identification based on electrocardiogram. \textit{Electrotechnical University "LETI"}, Saint-Petersburg, 2005.
\bibitem{3} Cui, W., Wang, Z., Li, Y. (2021). ECG-based biometric recognition under exercise and rest situations. \textit{Biomedical Engineering Advances}, 2, 100008.
\bibitem{4} Biel L., Pettersson O., Philipson L., Wide P. (2001). ECG analysis: a new approach in human identification. \textit{IEEE Transactions on Instrumentation and Measurement}, 50(3), 808-812.
\bibitem{5} Shen T.W., Tompkins W.J., Hu Y.H. (2002). One-lead ECG for identity verification. \textit{IEEE Engineering in Medicine and Biology Society}.
\bibitem{6} Analog Devices. \textit{AD620 Low Cost Low Power Instrumentation Amplifier Datasheet}.
\bibitem{7} Jaber, H. A., \& \c{C}ankaya, I. (2014). Heart Rate Monitoring and PQRST Detection Based on Graphical User Interface with Matlab. \textit{Yildirim Beyazit University}.
\bibitem{8} Madona, P., Basti, R. I., Zain, M. M. (2021). PQRST wave detection on ECG signals. \textit{Gac Sanit.}, 35(S2), S364-S369.
\end{thebibliography}

\end{document}